\documentclass[../../main.tex]{subfiles}% 注意这里的文件路径不能用 ./main.tex ,否则用latexmk编译子文件会报错
\graphicspath{{\subfix{./image/}}{\subfix{../../image/}}} % 指定图片引用路径,后续可以直接使用图片文件名
% 如果图片存放在上级目录,则使用 ../../image/ 作为备用路径

% 例如:
% \begin{figure}[H]
% \centering
% \includegraphics[width=0.3\textwidth]{图.png}
% \caption{}
% \label{figure:图}
% \end{figure}
% 注意:上述\label{}一定要放在\caption{}之后,否则引用图片序号会只会显示??.

\begin{document}
% SOURCE p=128; printed=113; object=ch04-topological-spaces-title; region=lower
\section{拓扑空间}\label{section:ch04-topological-spaces}
% SOURCE p=128; printed=113; object=ch04-topology-definition; region=lower
\begin{definition}[拓扑与拓扑空间]\label{definition:ch04-topology}
设 $X$ 为非空集合。$X$ 上的一个\textbf{拓扑}是一个 $X$ 的子集族 $\mathcal T$，它包含 $\varnothing$ 和 $X$，并且对任意并及有限交封闭：若 $\{U_\alpha\}_{\alpha\in A}\subseteq\mathcal T$，则 $\bigcup_{\alpha\in A}U_\alpha\in\mathcal T$；若 $U_1,\ldots,U_n\in\mathcal T$，则 $\bigcap_{j=1}^nU_j\in\mathcal T$。二元组 $(X,\mathcal T)$ 称为\textbf{拓扑空间}。若 $\mathcal T$ 已明确，就简称为拓扑空间 $X$。
\end{definition}
% SUPPLEMENT: id=v14-ch04-topology-remark; parent=ch04-topology-definition; kind=motivation+hypotheses; origin=independent-derivation
\begin{remark}
\textbf{动机与逐条件说明（重构新增）。}拓扑保留了“每个点附近有一小块空间”的信息，而不要求距离的数值。$\mathcal T\subseteq\mathscr P(X)$ 固定所有开集的共同底空间；$X$ 开使每个点都有邻域，$\varnothing$ 开处理没有点的情形。任意并封闭允许把任意多处的局部信息拼接成开集；有限交封闭允许同时满足有限个局部要求。一般不能要求任意交仍开：通常实直线中 $\bigcap_{n\ge1}(-1/n,1/n)=\{0\}$ 不是开集。若允许空指标族，则空并与空交分别为 $\varnothing$ 与 $X$，因而可把前两个要求并入后两个公理。非空性是本书的约定；同样的公理也容许空空间，其唯一拓扑为 $\{\varnothing\}$。后文遇到空子空间时，可采用这一自然延伸。
\end{remark}
% END SUPPLEMENT: id=v14-ch04-topology-remark
% SOURCE p=128; printed=113; object=ch04-topology-examples-intro; region=lower
来看几个例子。
% SOURCE p=128; printed=113; object=ch04-discrete-indiscrete; region=bottom
\begin{definition}[离散拓扑与平凡拓扑]\label{definition:ch04-discrete-indiscrete}
若 $X$ 为任意非空集合，则 $\mathscr P(X)$ 和 $\{\varnothing,X\}$ 都是 $X$ 上的拓扑，分别称为\textbf{离散拓扑}和\textbf{平凡拓扑}（或\textbf{不可分拓扑}）。
\end{definition}
% SUPPLEMENT: id=v14-ch04-discrete-indiscrete-details; parent=ch04-discrete-indiscrete; kind=proof+motivation+hypotheses; origin=independent-derivation
\begin{proof}
\textbf{补写验证（重构新增）。}$X$ 的子集的任意并及有限交仍为 $X$ 的子集，故 $\mathscr P(X)$ 满足全部公理。对 $\{\varnothing,X\}$，任意并只可能是 $\varnothing$ 或 $X$，有限交亦然，且两端集合都已包含。
\end{proof}
\begin{remark}
\textbf{动机与条件说明（重构新增）。}离散拓扑允许逐点区分，平凡拓扑只区分“全空间”和“没有点”，构成开集信息的两个极端。非空性沿用底空间约定；当 $X$ 只有一点时这两个拓扑相同，至少两点时才不同。
\end{remark}
% END SUPPLEMENT: id=v14-ch04-discrete-indiscrete-details
% SOURCE p=128; printed=113; object=ch04-cofinite-topology; region=bottom
\begin{definition}[余有限拓扑]\label{definition:ch04-cofinite-topology}
若 $X$ 为无限集合，则 $\{U\subseteq X:U=\varnothing\text{ 或 }U^c\text{ 有限}\}$ 是 $X$ 上的一个拓扑，称为\textbf{余有限拓扑}。
\end{definition}
% SUPPLEMENT: id=v14-ch04-cofinite-details; parent=ch04-cofinite-topology; kind=proof+motivation+hypotheses; origin=independent-derivation
\begin{proof}
\textbf{补写验证（重构新增）。}$\varnothing$ 已列入，且 $X^c=\varnothing$ 有限。任意开集族的并若非空，其中有一个非空成员 $U$；并集之补包含于有限集 $U^c$，故有限。有限个非空开集的交之补为有限个有限集的并，仍有限；若其中有空集，其交为空。故拓扑公理成立。
\end{proof}
\begin{remark}
\textbf{动机与条件说明（重构新增）。}“补集有限”表示开集仅排除有限个例外点。必须另列 $\varnothing$，因为无限 $X$ 的补集条件本身不包含空集。无限性并非公理验证所需；若 $X$ 有限，此构造等于离散拓扑，书中限制无限是为了得到不同于该退化情形的例子。
\end{remark}
% END SUPPLEMENT: id=v14-ch04-cofinite-details
% SOURCE p=129; printed=114; object=ch04-metric-topology; region=top
\begin{example}[度量诱导的拓扑]\label{example:ch04-metric-topology}
若 $X$ 是度量空间，则关于该度量的全体开集构成 $X$ 上的一个拓扑。
\end{example}
% SUPPLEMENT: id=v14-ch04-metric-topology-proof; parent=ch04-metric-topology; kind=proof+motivation+hypotheses; origin=independent-derivation
\begin{proof}
\textbf{补写验证（重构新增）。}空集的开性为空真，全空间包含每个点的任意开球。若点 $x$ 在开集族的并中，它属于某个成员，故有以 $x$ 为中心的开球包含于该成员及并集。若 $x$ 在有限个开集 $U_1,\ldots,U_n$ 的交中，对每个 $j$ 取 $r_j>0$ 使 $B(x,r_j)\subseteq U_j$；正数 $r=\min_jr_j$ 给出 $B(x,r)\subseteq\bigcap_jU_j$。
\end{proof}
\begin{remark}
\textbf{动机与条件说明（重构新增）。}此例保证抽象拓扑包含经典度量分析。有限性用于保证有限个正半径的最小值仍正；无限多个正半径的下确界可能为零，正是不能普遍取任意交的原因。
\end{remark}
% END SUPPLEMENT: id=v14-ch04-metric-topology-proof
% SOURCE p=129; printed=114; object=ch04-relative-topology; region=upper
\begin{definition}[相对拓扑]\label{definition:ch04-relative-topology}
若 $(X,\mathcal T)$ 为拓扑空间且 $Y\subseteq X$，则 $\mathcal T_Y=\{U\cap Y:U\in\mathcal T\}$ 是 $Y$ 上的拓扑，称为由 $\mathcal T$ 诱导的\textbf{相对拓扑}。
\end{definition}
% SUPPLEMENT: id=v14-ch04-relative-topology-details; parent=ch04-relative-topology; kind=proof+motivation+hypotheses; origin=independent-derivation
\begin{proof}
\textbf{补写验证（重构新增）。}$\varnothing=\varnothing\cap Y$，$Y=X\cap Y$。集合恒等式
\begin{align*}
\bigcup_\alpha(U_\alpha\cap Y)&=\left(\bigcup_\alpha U_\alpha\right)\cap Y, &
\bigcap_{j=1}^n(U_j\cap Y)&=\left(\bigcap_{j=1}^nU_j\right)\cap Y
\end{align*}
把所需封闭性化为 $\mathcal T$ 的封闭性。
\end{proof}
\begin{remark}
\textbf{动机与条件说明（重构新增）。}相对拓扑表示只在 $Y$ 内观察原空间的局部结构。要求 $Y\subseteq X$ 是为了作此限制；无需 $Y$ 在 $X$ 中开或闭。例如 $[0,1/2)$ 在 $Y=[0,1]$ 中相对开，却不在 $\mathbb R$ 中开。取 $Y=\varnothing$ 时采用前述空空间约定；若坚持空间非空，则相对拓扑的术语只对非空 $Y$ 使用。
\end{remark}
% END SUPPLEMENT: id=v14-ch04-relative-topology-details
% SOURCE p=129; printed=114; object=ch04-basic-terminology-intro; region=upper
现在介绍拓扑空间的基本术语，其中多数概念在度量空间中已经熟悉。在另行说明之前，$(X,\mathcal T)$ 始终是一个固定的拓扑空间。
% SOURCE p=129; printed=114; object=ch04-open-closed; region=upper
\begin{definition}[开集、闭集及相对开闭集]\label{definition:ch04-open-closed}
$\mathcal T$ 的成员称为\textbf{开集}，它们的补集称为\textbf{闭集}。若 $Y\subseteq X$，则 $Y$ 的子集在相对拓扑中开（闭）时，称为\textbf{相对开集}（\textbf{相对闭集}）。
\end{definition}
% SOURCE p=129; printed=114; object=ch04-closed-set-operations; region=upper
\begin{proposition}\label{proposition:ch04-closed-set-operations}
由 de Morgan 律，闭集族对任意交及有限并封闭。
\end{proposition}
% SUPPLEMENT: id=v14-ch04-open-closed-details; parent=ch04-open-closed; kind=proof+motivation+hypotheses; origin=independent-derivation
\begin{proof}
\textbf{闭集运算的补写证明（重构新增）。}若各 $F_\alpha$ 闭，则 $(\bigcap_\alpha F_\alpha)^c=\bigcup_\alpha F_\alpha^c$ 开；有限闭集并的补为有限个开集的交，也开。空交为 $X$，空并为空集，均闭。
\end{proof}
\begin{remark}
\textbf{动机与条件说明（重构新增）。}闭集与开集描述同一拓扑的互补信息；“闭”不是“不开”，因为一个集合可以同时开闭，也可以二者皆非。补集必须相对于正在讨论的底空间取得；$F\subseteq Y$ 相对闭，等价于 $F=Y\cap C$，其中 $C$ 在 $X$ 中闭：把 $Y\setminus F=Y\cap U$ 取 $Y$ 内的补集即可。任意交与有限并的区别直接来自开集公理，不能任意互换。
\end{remark}
% END SUPPLEMENT: id=v14-ch04-open-closed-details
% SOURCE p=129; printed=114; object=ch04-interior-closure; region=middle
\begin{definition}[内部与闭包]\label{definition:ch04-interior-closure}
若 $A\subseteq X$，包含在 $A$ 中的所有开集之并称为 $A$ 的\textbf{内部}，包含 $A$ 的所有闭集之交称为 $A$ 的\textbf{闭包}，分别记作 $A^\circ$ 和 $\overline A$。
\end{definition}
% SOURCE p=129; printed=114; object=ch04-interior-closure-properties; region=middle
\begin{proposition}\label{proposition:ch04-interior-closure-properties}
$A^\circ$ 是包含在 $A$ 中的最大开集，$\overline A$ 是包含 $A$ 的最小闭集，而且
\begin{align*}
(A^\circ)^c=\overline{A^c},\qquad (\overline A)^c=(A^c)^\circ.
\end{align*}
\end{proposition}
% SUPPLEMENT: id=v14-ch04-interior-closure-details; parent=ch04-interior-closure-properties; kind=proof+motivation+hypotheses; origin=independent-derivation
\begin{proof}
\textbf{补写证明（重构新增）。}定义中的并仍开，且包含在 $A$ 中；任何包含于 $A$ 的开集都已列入此并，故为最大。闭包的最小性同理由任意闭集交仍闭得到。条件 $U\subseteq A$ 与 $A^c\subseteq U^c$ 等价，而开集 $U$ 与闭集 $U^c$ 一一对应；将定义中的并取补，de Morgan 律便给出第一式。以 $A^c$ 替换 $A$ 再取补得到第二式。
\end{proof}
\begin{remark}
\textbf{动机与条件说明（重构新增）。}内部从集合内侧提取开集，闭包从外侧补成闭集；取所有候选集使结果唯一且不依赖选取。候选族不会造成存在性障碍：$\varnothing\subseteq A\subseteq X$。不要求 $A$ 本身开、闭或可测，也不要求任何分离公理。两式说明这两种操作由补集互相决定。
\end{remark}
% END SUPPLEMENT: id=v14-ch04-interior-closure-details
% SOURCE p=129; printed=114; object=ch04-boundary-density; region=middle
\begin{definition}[边界、稠密与无处稠密]\label{definition:ch04-boundary-density}
差集 $\overline A\setminus A^\circ=\overline A\cap\overline{A^c}$ 称为 $A$ 的\textbf{边界}，记作 $\partial A$。若 $\overline A=X$，称 $A$ 在 $X$ 中\textbf{稠密}。若 $(\overline A)^\circ=\varnothing$，称 $A$ \textbf{无处稠密}。
\end{definition}
% SUPPLEMENT: id=v14-ch04-boundary-density-details; parent=ch04-boundary-density; kind=proof+motivation+hypotheses; origin=independent-derivation
\begin{remark}
\textbf{动机与逐条件说明（重构新增）。}边界由既不能从开邻域中排除 $A$、也不能排除 $A^c$ 的点组成；公式来自 $(A^\circ)^c=\overline{A^c}$。稠密的闭包等于整个底空间，表示每个非空开集都遇到 $A$：否则某个非空开集与 $A$ 不交，其补是包含 $A$ 的真闭集；反向若闭包不满，$(\overline A)^c$ 就是这样的开集。无处稠密要求先取闭包再取内部，仅有 $A^\circ=\varnothing$ 不够，例如 $\mathbb Q$ 在 $\mathbb R$ 中内部为空却稠密。等价地，每个非空开集 $U$ 都包含一个与 $A$ 不交的非空开集：前向取 $U\setminus\overline A$；反向若 $U\subseteq\overline A$ 非空开，则 $U$ 中任意非空开子集都遇到 $A$，矛盾。
\end{remark}
% END SUPPLEMENT: id=v14-ch04-boundary-density-details
% SOURCE p=129; printed=114; object=ch04-neighborhood; region=middle
\begin{definition}[邻域]\label{definition:ch04-neighborhood}
若 $x\in X$（或 $E\subseteq X$），$x$（或 $E$）的一个\textbf{邻域}是满足 $x\in A^\circ$（或 $E\subseteq A^\circ$）的集合 $A\subseteq X$。因此 $A$ 开当且仅当它是自身的邻域。（有些作者要求邻域是开集，本书不作此要求。）
\end{definition}
% SUPPLEMENT: id=v14-ch04-neighborhood-details; parent=ch04-neighborhood; kind=motivation+hypotheses; origin=independent-derivation
\begin{remark}
\textbf{动机与条件说明（重构新增）。}使用内部而非仅要求 $x\in A$，是为了使 $A$ 包含围绕 $x$ 的某个开集；例如 $\{0\}$ 含有 $0$，却不是通常实直线中 $0$ 的邻域。允许 $A$ 不开使闭球等集合也可充当邻域。集合 $E$ 的邻域须同时包含其每一点附近的开集，但不要求存在某个统一半径。$A$ 是自身邻域等价于 $A\subseteq A^\circ$；结合反向包含即得 $A=A^\circ$，也即 $A$ 开。特别地，不可从“$V$ 是邻域”直接推断“$V$ 开”。
\end{remark}
% END SUPPLEMENT: id=v14-ch04-neighborhood-details
% SOURCE p=129; printed=114; object=ch04-accumulation-point; region=middle
\begin{definition}[聚点]\label{definition:ch04-accumulation-point}
若对 $x$ 的每个邻域 $U$ 都有 $A\cap(U\setminus\{x\})\ne\varnothing$，则称 $x$ 为 $A$ 的\textbf{聚点}。其他著作有时用“簇点”（cluster point）或“极限点”（limit point）表示同一概念；下文将把“簇点”用于一个略有不同的概念。
\end{definition}
% SUPPLEMENT: id=v14-ch04-accumulation-remark; parent=ch04-accumulation-point; kind=motivation+hypotheses; origin=independent-derivation
\begin{remark}
\textbf{动机与条件说明（重构新增）。}删去 $\{x\}$ 是为了排除由点自身造成的平凡接触；否则每个 $x\in A$ 都自动满足。要求所有邻域都遇到其余点，表示这种接触不能用缩小邻域消除。只检查开邻域已经充分，因为每个邻域包含开邻域。不预设 $T_1$ 时，不能把“遇到一个异于 $x$ 的点”擅自换成“遇到无限多个点”：在两点平凡空间中，$A$ 为单点时，另一点就是 $A$ 的聚点。
\end{remark}
% END SUPPLEMENT: id=v14-ch04-accumulation-remark
% SOURCE p=129; printed=114; object=ch04-closure-accumulation; number=4.1; region=lower-middle
\begin{proposition}\label{proposition:ch04-closure-accumulation}
设 $A\subseteq X$，用 $\operatorname{acc}(A)$ 表示 $A$ 的聚点集，则 $\overline A=A\cup\operatorname{acc}(A)$，并且 $A$ 闭当且仅当 $\operatorname{acc}(A)\subseteq A$。
\end{proposition}
% SOURCE p=129; printed=114; object=ch04-proof-closure-accumulation; region=lower-middle
\begin{proof}
若 $x\notin\overline A$，则 $A^c$ 是 $x$ 的一个与 $A$ 不交的邻域，故 $x\notin\operatorname{acc}(A)$；因此 $A\cup\operatorname{acc}(A)\subseteq\overline A$。若 $x\notin A\cup\operatorname{acc}(A)$，则存在包含 $x$ 的开集 $U$，使 $U\cap A=\varnothing$，于是 $\overline A\subseteq U^c$ 且 $x\notin\overline A$。所以 $\overline A\subseteq A\cup\operatorname{acc}(A)$。最后，$A$ 闭当且仅当 $A=\overline A$，而这等价于 $\operatorname{acc}(A)\subseteq A$。
\end{proof}
% SUPPLEMENT: id=v14-ch04-closure-accumulation-remark; parent=ch04-closure-accumulation; kind=proof-detail+motivation+hypotheses; origin=independent-derivation
\begin{remark}
\textbf{动机、条件与证明细节（重构新增）。}该命题把闭包分成原集合的点和必须补入的聚点；并集中的 $A$ 不能删除，例如离散空间中非空 $A$ 没有聚点而闭包仍是 $A$。证明首句中，$(\overline A)^c$ 是包含 $x$ 且包含于 $A^c$ 的开集，故 $A^c$ 确为邻域。第二方向由非聚点得到某邻域 $N$ 满足 $A\cap(N\setminus\{x\})=\varnothing$，再由 $x\notin A$ 得 $A\cap N=\varnothing$，取 $U=N^\circ$ 即可。证明完全不需要度量、可数性或分离公理。
\end{remark}
% END SUPPLEMENT: id=v14-ch04-closure-accumulation-remark
% SOURCE p=129; printed=114; object=ch04-weaker-stronger; region=lower
\begin{definition}[较弱与较强的拓扑]\label{definition:ch04-weaker-stronger}
若 $\mathcal T_1,\mathcal T_2$ 是 $X$ 上的拓扑且 $\mathcal T_1\subseteq\mathcal T_2$，则称 $\mathcal T_1$ 比 $\mathcal T_2$ \textbf{弱}（或\textbf{粗}），或称 $\mathcal T_2$ 比 $\mathcal T_1$ \textbf{强}（或\textbf{细}）。显然，平凡拓扑是 $X$ 上最弱的拓扑，离散拓扑是最强的拓扑。
\end{definition}
% SUPPLEMENT: id=v14-ch04-weaker-stronger-remark; parent=ch04-weaker-stronger; kind=motivation+hypotheses; origin=independent-derivation
\begin{remark}
\textbf{动机与条件说明（重构新增）。}在同一集合 $X$ 上比较开集数量，才能比较局部信息的精细程度；开集越多，邻域测试越多。包含关系允许相等，不表示严格强弱；两个拓扑也可能不可比较。每个拓扑都含 $\varnothing,X$ 且包含于 $\mathscr P(X)$，直接给出两个极端的断言。
\end{remark}
% END SUPPLEMENT: id=v14-ch04-weaker-stronger-remark
% SOURCE p=129; printed=114; object=ch04-generated-topology; region=lower
\begin{definition}[生成拓扑与子基]\label{definition:ch04-generated-topology}
若 $\mathcal E\subseteq\mathscr P(X)$，则有唯一最弱的拓扑 $\mathcal T(\mathcal E)$ 包含 $\mathcal E$，即所有包含 $\mathcal E$ 的 $X$ 上拓扑之交。它称为由 $\mathcal E$ \textbf{生成的拓扑}；有时称 $\mathcal E$ 为 $\mathcal T(\mathcal E)$ 的一个\textbf{子基}。
\end{definition}
% SUPPLEMENT: id=v14-ch04-generated-topology-details; parent=ch04-generated-topology; kind=proof+motivation+hypotheses; origin=independent-derivation
\begin{proof}
\textbf{存在性与唯一性的补写证明（重构新增）。}候选拓扑族非空，因为离散拓扑包含 $\mathcal E$。所有候选族的交含 $\varnothing,X,\mathcal E$；其成员的任意并及有限交属于每一个候选拓扑，故也属于交。所得族为拓扑，并包含于任何候选拓扑，因而最弱。若另一个也最弱，则二者互相包含，所以相等。
\end{proof}
\begin{remark}
\textbf{动机与条件说明（重构新增）。}生成拓扑只强制指定集合开，再补入公理不可避免的开集。$\mathcal E$ 只需由 $X$ 的子集组成，不必覆盖 $X$，也不必对任何集合运算封闭，甚至可以为空；当它为空时得到平凡拓扑。这里取的是“所有含 $\mathcal E$ 的拓扑”的交，并非直接取 $\mathcal E$ 中各集合的交。
\end{remark}
% END SUPPLEMENT: id=v14-ch04-generated-topology-details
% SOURCE p=129; printed=114; object=ch04-neighborhood-base; region=bottom
\begin{definition}[邻域基]\label{definition:ch04-neighborhood-base}
设 $\mathcal T$ 是 $X$ 上的拓扑。$\mathcal T$ 在 $x\in X$ 处的一个\textbf{邻域基}是集合族 $\mathcal N\subseteq\mathcal T$，满足：
\begin{enumerate}
\item 对所有 $V\in\mathcal N$，有 $x\in V$；
\item 若 $U\in\mathcal T$ 且 $x\in U$，则存在 $V\in\mathcal N$ 使 $V\subseteq U$。
\end{enumerate}
\end{definition}
% SUPPLEMENT: id=v14-ch04-neighborhood-base-remark; parent=ch04-neighborhood-base; kind=motivation+hypotheses; origin=independent-derivation
\begin{remark}
\textbf{动机与逐条件说明（重构新增）。}邻域基是一族足以代替全部邻域作局部测试的集合。本书明确要求成员开；含点条件使每个成员都是 $x$ 的开邻域；细化条件保证没有遗漏任何局部要求。将 $U=X$ 代入可知邻域基不为空。虽然一般邻域允许不开放，任何一般邻域都包含一个开邻域，因而也包含某个基成员；这解释了定义只测试开 $U$ 的充分性。
\end{remark}
% END SUPPLEMENT: id=v14-ch04-neighborhood-base-remark
% SOURCE p=130; printed=115; object=ch04-base-definition; region=top
\begin{definition}[拓扑的基]\label{definition:ch04-topology-base}
$\mathcal T$ 的一个\textbf{基}是集合族 $\mathcal B\subseteq\mathcal T$，它在每个 $x\in X$ 处都包含 $\mathcal T$ 的一个邻域基。
\end{definition}
% SOURCE p=130; printed=115; object=ch04-metric-bases; region=top
\begin{example}\label{example:ch04-metric-bases}
若 $X$ 为度量空间，则以 $x$ 为中心的全体开球是度量拓扑在 $x$ 处的邻域基，而 $X$ 中所有开球组成一个基。
\end{example}
% SUPPLEMENT: id=v14-ch04-base-metric-details; parent=ch04-base-definition; kind=proof+motivation+hypotheses; origin=independent-derivation
\begin{proof}
\textbf{度量情形的补写验证（重构新增）。}对开集 $U\ni x$，开性的定义给出 $r>0$ 使 $B(x,r)\subseteq U$。球本身开且含其中心，故满足邻域基的两个要求。所有开球的族包含每个点的这族中心球，故为基。
\end{proof}
\begin{remark}
\textbf{动机与条件说明（重构新增）。}邻域基控制一个固定点，拓扑基同时控制所有点；要求基成员开，才能用其并恢复原拓扑。基不必由两两不同形状的集合组成，也不必可数；可数性将在后面另作公理。定义只要求“包含”各点的邻域基，并非每个基成员都要含每一个点。
\end{remark}
% END SUPPLEMENT: id=v14-ch04-base-metric-details
% SOURCE p=130; printed=115; object=ch04-base-union-characterization; number=4.2; region=upper
\begin{proposition}\label{proposition:ch04-base-union-characterization}
若 $\mathcal T$ 是 $X$ 上的拓扑，且 $\mathcal E\subseteq\mathcal T$，则 $\mathcal E$ 是 $\mathcal T$ 的一个基，当且仅当每个非空 $U\in\mathcal T$ 都是 $\mathcal E$ 中一些成员的并。
\end{proposition}
% SOURCE p=130; printed=115; object=ch04-proof-base-union-characterization; region=upper
\begin{proof}
若 $\mathcal E$ 是基，$U\in\mathcal T$ 且 $x\in U$，则存在 $V_x\in\mathcal E$ 使 $x\in V_x\subseteq U$，故 $U=\bigcup_{x\in U}V_x$。反之，若每个非空 $U\in\mathcal T$ 都是 $\mathcal E$ 中成员的并，则 $\{V\in\mathcal E:x\in V\}$ 显然是 $x$ 处的邻域基，所以 $\mathcal E$ 为基。
\end{proof}
% SUPPLEMENT: id=v14-ch04-base-union-remark; parent=ch04-base-union-characterization; kind=proof-detail+motivation+hypotheses; origin=independent-derivation
\begin{remark}
\textbf{动机、条件与证明细节（重构新增）。}此判别把逐点细化改写成开集的整体分解。$\mathcal E\subseteq\mathcal T$ 保证候选成员确实开；若删去，可让单点族把任何集合都写成并，却通常不是原拓扑的基。非空的限制仅为避免讨论空并，约定空并为空集后可去掉。逆向中，对任意 $x\in U$，从表示 $U=\bigcup_\alpha V_\alpha$ 选一个含 $x$ 的成员，它自然满足 $V_\alpha\subseteq U$，故邻域基的细化条件成立。
\end{remark}
% END SUPPLEMENT: id=v14-ch04-base-union-remark
% SOURCE p=130; printed=115; object=ch04-base-criterion; number=4.3; region=middle
\begin{proposition}\label{proposition:ch04-base-criterion}
若 $\mathcal E\subseteq\mathscr P(X)$，则 $\mathcal E$ 是 $X$ 上某个拓扑的基，当且仅当满足以下两个条件：
\begin{enumerate}
\item\label{item:ch04-base-cover} 每个 $x\in X$ 都属于某个 $V\in\mathcal E$；
\item\label{item:ch04-base-intersection} 若 $U,V\in\mathcal E$ 且 $x\in U\cap V$，则存在 $W\in\mathcal E$ 使 $x\in W\subseteq U\cap V$。
\end{enumerate}
\end{proposition}
% SOURCE p=130; printed=115; object=ch04-proof-base-criterion; region=middle
\begin{proof}
必要性显然，因为若 $U,V$ 开，则 $U\cap V$ 也开。为证充分性，令
\begin{align*}
\mathcal T=\{U\subseteq X:\text{对每个 }x\in U,\text{存在 }V\in\mathcal E\text{ 使 }x\in V\subseteq U\}.
\end{align*}
由\subcref{proposition:ch04-base-criterion}{item:ch04-base-cover}有 $X\in\mathcal T$，显然 $\varnothing\in\mathcal T$，且 $\mathcal T$ 显然对任意并封闭。若 $U_1,U_2\in\mathcal T$ 且 $x\in U_1\cap U_2$，则存在 $V_1,V_2\in\mathcal E$，使 $x\in V_1\subseteq U_1$ 及 $x\in V_2\subseteq U_2$；再由\subcref{proposition:ch04-base-criterion}{item:ch04-base-intersection}，有 $W\in\mathcal E$ 使 $x\in W\subseteq V_1\cap V_2$。于是 $U_1\cap U_2\in\mathcal T$，由归纳法，$\mathcal T$ 对有限交封闭。因此 $\mathcal T$ 为拓扑，而 $\mathcal E$ 显然是其一个基。
\end{proof}
% SUPPLEMENT: id=v14-ch04-base-criterion-details; parent=ch04-base-criterion; kind=proof-detail+motivation+hypotheses; origin=independent-derivation
\begin{remark}
\textbf{动机、逐条件说明与省略步骤（重构新增）。}覆盖条件对应全空间开；交的局部细化条件对应有限交开，却不要求 $U\cap V$ 本身属于基。必要性中，用基在 $X$ 中逐点细化得\subcref{proposition:ch04-base-criterion}{item:ch04-base-cover}，在开集 $U\cap V$ 中逐点细化得\subcref{proposition:ch04-base-criterion}{item:ch04-base-intersection}。充分性中的并封闭为：点在某个 $U_\alpha$ 中时，选包含于 $U_\alpha$ 的基成员，它也包含于并集。任意 $V\in\mathcal E$ 自己就为其每一点提供所需基成员，故 $\mathcal E\subseteq\mathcal T$；定义中的局部细化条件随后证明它是基。

两条件不能互相替代：在 $X=\{1,2\}$ 上取 $\mathcal E=\{\{1\}\}$，满足交条件却不覆盖；在 $X=\{1,2,3\}$ 上取 $\mathcal E=\{\{1,2\},\{2,3\}\}$，覆盖全空间，却不能在两成员的交 $\{2\}$ 内作所需细化。所成拓扑唯一，因为上一命题已确定其成员恰为基成员的任意并。
\end{remark}
% END SUPPLEMENT: id=v14-ch04-base-criterion-details
% SOURCE p=130; printed=115; object=ch04-subbase-construction; number=4.4; region=lower
\begin{proposition}\label{proposition:ch04-subbase-construction}
若 $\mathcal E\subseteq\mathscr P(X)$，则由 $\mathcal E$ 生成的拓扑 $\mathcal T(\mathcal E)$ 由 $\varnothing,X$ 以及 $\mathcal E$ 的成员的有限交的所有并组成。
\end{proposition}
% SOURCE p=130; printed=115; object=ch04-proof-subbase-construction; region=lower
\begin{proof}
$\mathcal E$ 中集合的有限交所成的族，加上 $X$，满足\cref{proposition:ch04-base-criterion}的条件。因此，由\cref{proposition:ch04-base-union-characterization}，这类集合的所有并再加上 $\varnothing$，组成一个拓扑。它显然包含于 $\mathcal T(\mathcal E)$，所以等于 $\mathcal T(\mathcal E)$。
\end{proof}
% SUPPLEMENT: id=v14-ch04-subbase-construction-details; parent=ch04-subbase-construction; kind=proof-detail+motivation+hypotheses; origin=independent-derivation
\begin{remark}
\textbf{动机、条件与最小性细节（重构新增）。}构造按“先有限交、后任意并”的顺序补齐两项公理。将空交约定为 $X$、空并约定为 $\varnothing$ 后，无需另列两端集合；原文单独列出它们兼容不采用空运算的写法。有限交族包含 $X$，因而覆盖；两个有限交再相交仍为有限交，故满足基判别。所得拓扑 $\mathcal S$ 含每个 $E\in\mathcal E$（取单项交），所以由生成拓扑的最小性有 $\mathcal T(\mathcal E)\subseteq\mathcal S$；另一方面每一构造步骤都保留在 $\mathcal T(\mathcal E)$ 内，给出反向包含。这里不要求子基覆盖空间，$\mathcal E=\varnothing$ 也得到 $\{\varnothing,X\}$。
\end{remark}
% END SUPPLEMENT: id=v14-ch04-subbase-construction-details
% SOURCE p=130; printed=115; object=ch04-subbase-comparison; region=lower
注意，这一命题的简洁程度与 $\sigma$-代数的相应结果（\cref{proposition:ch01-transfinite-sigma-generation}）形成对比。这里之所以更容易，是因为只涉及有限交。
% SOURCE p=130; printed=115; object=ch04-axioms-introduction; region=bottom
拓扑空间的概念足够一般，能够涵盖大量有趣的例子；但同样因为过于一般，无法直接得到很多有趣的定理。为了建立有用的理论，通常必须限制所考虑的空间类别。本节余下部分将讨论两类常见限制，即所谓可数性公理与分离公理。
% SOURCE p=131; printed=116; object=ch04-first-countability; region=top
\begin{definition}[第一可数性]\label{definition:ch04-first-countable}
若拓扑空间 $(X,\mathcal T)$ 的每一点都有可数邻域基，则称它满足\textbf{第一可数性公理}，或称它是\textbf{第一可数的}。
\end{definition}
% SOURCE p=131; printed=116; object=ch04-nested-countable-base; region=top
\begin{proposition}\label{proposition:ch04-nested-countable-base}
若 $X$ 第一可数，则对每个 $x\in X$，存在 $x$ 处的邻域基 $\{U_j\}_{j=1}^\infty$，满足对所有 $j$ 有 $U_j\supseteq U_{j+1}$。事实上，若 $\{V_j\}_{j=1}^\infty$ 是任一可数邻域基，可以取 $U_j=\bigcap_{i=1}^jV_i$。
\end{proposition}
% SUPPLEMENT: id=v14-ch04-first-countable-details; parent=ch04-first-countability; kind=proof+motivation+hypotheses; origin=independent-derivation
\begin{proof}
\textbf{递减邻域基的补写验证（重构新增）。}每个 $U_j$ 是有限个开邻域的交，故开且含 $x$；显然逐项递减。若 $W$ 是 $x$ 的开邻域，选 $V_k\subseteq W$，则 $U_k\subseteq V_k\subseteq W$，所以新族仍为邻域基。有限邻域基可重复枚举成序列，不影响论证。
\end{proof}
\begin{remark}
\textbf{动机与条件说明（重构新增）。}第一可数性允许把一个点附近的全部局部测试压缩为一条序列；不要求全空间共用同一可数族。递减化利用有限交，避免错误地假设无限交仍是邻域。每个点都满足是空间的整体公理，但后面在某个固定点用序列描述闭包，只需该点有可数邻域基。
\end{remark}
% END SUPPLEMENT: id=v14-ch04-first-countable-details
% SOURCE p=131; printed=116; object=ch04-second-countability; region=top
\begin{definition}[第二可数性]\label{definition:ch04-second-countable}
若 $\mathcal T$ 有一个可数基，则称 $(X,\mathcal T)$ 满足\textbf{第二可数性公理}，或称它是\textbf{第二可数的}。
\end{definition}
% SUPPLEMENT: id=v14-ch04-second-countable-remark; parent=ch04-second-countability; kind=motivation+hypotheses; origin=independent-derivation
\begin{remark}
\textbf{动机与条件说明（重构新增）。}第二可数性要求一族统一的可数开集服务于所有点，比逐点各有一族更强。若 $\mathcal B$ 可数，则 $\{B\in\mathcal B:x\in B\}$ 是可数邻域基，所以第二可数推出第一可数。可数的是基的成员数，并不意味着所有开集构成可数集合，也不意味着底空间可数。
\end{remark}
% END SUPPLEMENT: id=v14-ch04-second-countable-remark
% SOURCE p=131; printed=116; object=ch04-separability; region=top
\begin{definition}[可分空间]\label{definition:ch04-separable}
若 $X$ 有一个可数稠密子集，则称 $(X,\mathcal T)$ \textbf{可分}。
\end{definition}
% SUPPLEMENT: id=v14-ch04-separable-remark; parent=ch04-separability; kind=motivation+hypotheses; origin=independent-derivation
\begin{remark}
\textbf{动机与条件说明（重构新增）。}可分性要求可数个点组成的“采样集”遇到每个非空开集。稠密性不可删去，否则任意非空空间都有单点可数子集，条件就失去内容；只要求稠密而不要求可数也无约束，因为 $X$ 自身总稠密。这里“可分”与分离公理中的“分离”不是同一性质。
\end{remark}
% END SUPPLEMENT: id=v14-ch04-separable-remark
% SOURCE p=131; printed=116; object=ch04-metric-countability; region=upper
\begin{proposition}\label{proposition:ch04-metric-countability}
每个度量空间都第一可数（以 $x$ 为中心、正有理数为半径的球组成 $x$ 处的邻域基）；度量空间第二可数当且仅当它可分（\cref{exercise:ch04-separable-metric-second-countable}）。
\end{proposition}
% SUPPLEMENT: id=v14-ch04-metric-countability-details; parent=ch04-metric-countability; kind=proof+motivation+hypotheses; origin=independent-derivation
\begin{proof}
\textbf{补写证明（重构新增）。}给定 $B(x,r)$，选正有理数 $q<r$，即有 $B(x,q)\subseteq B(x,r)$；正有理数可数，故第一断言成立。可分推出第二可数的完整构造见\cref{exercise:ch04-separable-metric-second-countable}的解答，反向由下一命题给出；这两处证明均不依赖本断言。
\end{proof}
\begin{remark}
\textbf{动机与条件说明（重构新增）。}度量提供可数的半径刻度；可数稠密集进一步提供可数的中心选择。一般拓扑空间中没有这两种数值结构，所以可分性不能普遍推出第二可数性，\cref{exercise:ch04-upper-limit-topology}给出反例。
\end{remark}
% END SUPPLEMENT: id=v14-ch04-metric-countability-details
% SOURCE p=131; printed=116; object=ch04-second-countable-separable-intro; region=upper
后一事实可以部分推广：
% SOURCE p=131; printed=116; object=ch04-second-countable-separable; number=4.5; region=upper-middle
\begin{proposition}\label{proposition:ch04-second-countable-separable}
每个第二可数空间都可分。
\end{proposition}
% SOURCE p=131; printed=116; object=ch04-proof-second-countable-separable; region=upper-middle
\begin{proof}
若 $X$ 第二可数，取拓扑的一个可数基 $\mathcal E$，并对每个 $U\in\mathcal E$ 选一点 $x_U\in U$。则 $\{x_U:U\in\mathcal E\}$ 的闭包之补是一个不包含任何 $U\in\mathcal E$ 的开集，因而为空集；所以 $\{x_U:U\in\mathcal E\}$ 稠密。
\end{proof}
% SUPPLEMENT: id=v14-ch04-second-countable-separable-details; parent=ch04-second-countable-separable; kind=proof-correction+motivation+hypotheses; origin=independent-derivation
\begin{remark}
\textbf{动机、条件与空基成员的澄清（重构新增）。}原证明的“对每个 $U$ 选一点”隐含先删除可能存在的空集成员；定义并没有禁止基含空集。令 $\mathcal E'=\mathcal E\setminus\{\varnothing\}$，仍是可数基，再作选择。所得点集 $D$ 可数；每个非空开集 $V$ 含一个非空基成员 $U$，从而包含 $x_U\in D$，故 $D$ 稠密。可数基的作用是保证所选点集可数，不需要任何分离公理。逆命题在一般空间中不成立，见\cref{exercise:ch04-upper-limit-topology}。
\end{remark}
% END SUPPLEMENT: id=v14-ch04-second-countable-separable-details
% SOURCE p=131; printed=116; object=ch04-sequence-convergence; region=middle
\begin{definition}[序列收敛]\label{definition:ch04-sequence-convergence}
拓扑空间 $X$ 中的序列 $\{x_j\}$ \textbf{收敛}于 $x\in X$（记作 $x_j\to x$），是指对 $x$ 的每个邻域 $U$，存在 $J\in\mathbb N$，使所有 $j>J$ 都有 $x_j\in U$。
\end{definition}
% SUPPLEMENT: id=v14-ch04-sequence-convergence-remark; parent=ch04-sequence-convergence; kind=motivation+hypotheses; origin=independent-derivation
\begin{remark}
\textbf{动机与逐条件说明（重构新增）。}此定义把度量中的“最终距离小于任意正数”换成“最终进入任意邻域”。$J$ 可以依赖邻域，但其后的所有项必须进入；仅有无穷多项进入是簇点条件，较弱。只检查邻域基即可，因为每个邻域包含基成员。未附加 Hausdorff 条件时极限不必唯一，\cref{exercise:ch04-cofinite-sequence}将给出具体例子。
\end{remark}
% END SUPPLEMENT: id=v14-ch04-sequence-convergence-remark
% SOURCE p=131; printed=116; object=ch04-sequential-closure-intro; region=middle
第一可数空间有一个方便的性质：闭包、连续性等概念可以用序列收敛刻画；更一般的空间未必如此，后面将会看到。例如：
% SOURCE p=131; printed=116; object=ch04-sequential-closure; number=4.6; region=middle-lower
\begin{proposition}\label{proposition:ch04-sequential-closure}
若 $X$ 第一可数且 $A\subseteq X$，则 $x\in\overline A$ 当且仅当 $A$ 中存在一个收敛于 $x$ 的序列 $\{x_j\}$。
\end{proposition}
% SOURCE p=131; printed=116; object=ch04-proof-sequential-closure; region=middle-lower
\begin{proof}
取 $x$ 处的一个可数邻域基 $\{U_j\}$，满足所有 $j$ 都有 $U_j\supseteq U_{j+1}$。若 $x\in\overline A$，则对所有 $j$，有 $U_j\cap A\ne\varnothing$。选取 $x_j\in U_j\cap A$；当 $k>j$ 时 $U_k\subseteq U_j$，而 $x$ 的每个邻域都包含某个 $U_j$，所以显然 $x_j\to x$。另一方面，若 $x\notin\overline A$，且 $\{x_j\}$ 为 $A$ 中任意序列，则 $(\overline A)^c$ 是 $x$ 的邻域，却不含任何 $x_j$，故 $x_j\not\to x$。
\end{proof}
% SUPPLEMENT: id=v14-ch04-sequential-closure-remark; parent=ch04-sequential-closure; kind=motivation+hypotheses+counterexample; origin=independent-derivation
\begin{remark}
\textbf{动机与条件说明（重构新增）。}第一可数性用于把全部邻域压缩为递减序列，从而每一步只选一个点。反方向不需要第一可数性，正方向也只需在指定点 $x$ 处第一可数。$x\in A$ 时允许常值序列；$A=\varnothing$ 时两边都不成立。

一般不能删掉可数性控制。取不可数集合 $X$，令开集为 $\varnothing$ 及补集可数的集合；任意非空开集并的补包含于一个可数集，有限交的补为有限个可数集的并，故这确为拓扑。固定 $x$，$A=X\setminus\{x\}$ 稠密，因为任何非空开集都不可数，从而必遇到 $A$。但对 $A$ 中任意序列，集合 $U=X\setminus\{x_j:j\in\mathbb N\}$ 是 $x$ 的开邻域，且不含任何一项，故序列不能收敛于 $x$。这说明一般拓扑的闭包可能无法用序列检测。
\end{remark}
% END SUPPLEMENT: id=v14-ch04-sequential-closure-remark
% SOURCE p=131; printed=116; object=ch04-separation-intro; region=lower
最后讨论分离公理。这些拓扑空间的性质记作 $T_0,\ldots,T_4$，它们保证存在把点或闭集彼此分开的开集。若 $X$ 有性质 $T_j$，就称 $X$ 是 $T_j$ 空间，或称 $X$ 上的拓扑是 $T_j$ 的。
% SOURCE p=131,132; printed=116,117; object=ch04-separation-axioms; region=bottom-top
\begin{definition}[分离公理]\label{definition:ch04-separation-axioms}
\begin{description}
\item[$T_0$：] 若 $x\ne y$，则存在含 $x$ 不含 $y$ 的开集，或存在含 $y$ 不含 $x$ 的开集。
\item[$T_1$：] 若 $x\ne y$，则存在含 $y$ 不含 $x$ 的开集。
\item[$T_2$：] 若 $x\ne y$，则存在不交开集 $U,V$，使 $x\in U$ 且 $y\in V$。
\item[$T_3$：] $X$ 是 $T_1$ 空间，且对任意闭集 $A\subseteq X$ 及任意 $x\in A^c$，存在不交开集 $U,V$，使 $x\in U$ 且 $A\subseteq V$。
\item[$T_4$：] $X$ 是 $T_1$ 空间，且对任意不交闭集 $A,B\subseteq X$，存在不交开集 $U,V$，使 $A\subseteq U$ 且 $B\subseteq V$。
\end{description}
\end{definition}
% SOURCE p=132; printed=117; object=ch04-hausdorff-regular-normal; region=upper
\begin{definition}[Hausdorff 空间、正则空间与正规空间]\label{definition:ch04-hausdorff-regular-normal}
$T_2,T_3,T_4$ 还有其他名称：$T_2$ 空间称为 \textbf{Hausdorff 空间}，$T_3$ 空间称为\textbf{正则空间}，$T_4$ 空间称为\textbf{正规空间}。（有些作者不要求正则及正规空间为 $T_1$。）
\end{definition}
% SUPPLEMENT: id=v14-ch04-separation-remarks; parent=ch04-separation-axioms; kind=motivation+hypotheses; origin=independent-derivation
\begin{remark}
\textbf{动机与逐条件说明（重构新增）。}$T_0$ 要求至少能从一个方向区分两点；$T_1$ 的量词遍历所有有序对，因此交换 $x,y$ 后也可反向区分，但两次取得的开集未必不交。$T_2$ 要求同时给出互不相交的邻域，用于排除同一序列有两个不同极限：若序列最终进入 $U$，也最终进入 $V$，则足够后的项会在空集 $U\cap V$ 内，矛盾。$x\ne y$ 不可删去，否则同一点不可能落入不交开集。

$T_3$ 将点与点的分离提升为点与闭集的分离；$x\notin A$ 是可能分离的必要条件。闭性也有作用：即使在实直线上，$x=0$ 与 $A=\mathbb R\setminus\{0\}$ 不交，也无法用不交开集把它们包住，因为 $0$ 的每个开邻域都遇到 $A$。$T_4$ 进一步同时分离两个闭集；不交条件不可删去，闭性不能任意取消，例如稠密且不交的 $\mathbb Q$ 与其补无法如此分离。对空闭集，可取空开邻域，公理仍成立。

本书把 $T_1$ 明列在 $T_3,T_4$ 中，使单点闭性可以参与分离。若删去它，至少两点的平凡空间也满足这两项闭集分离条件，却不 Hausdorff：它仅有闭集 $\varnothing,X$，所需的非平凡分离根本不会触发。术语的这一约定在引用其他文献时应保持一致。
\end{remark}
% END SUPPLEMENT: id=v14-ch04-separation-remarks
% SOURCE p=132; printed=117; object=ch04-separation-hierarchy-intro; region=upper
还有一个介于 $T_3$ 与 $T_4$ 之间的有用分离条件，将在\cref{section:ch04-continuous-maps}讨论。下面对 $T_1$ 空间的刻画很有用，特别地，它说明每个正规空间都正则，每个正则空间都 Hausdorff。
% SOURCE p=132; printed=117; object=ch04-t1-singletons; number=4.7; region=upper-middle
\begin{proposition}\label{proposition:ch04-t1-singletons}
$X$ 是 $T_1$ 空间，当且仅当对每个 $x\in X$，$\{x\}$ 都是闭集。
\end{proposition}
% SOURCE p=132; printed=117; object=ch04-proof-t1-singletons; region=upper-middle
\begin{proof}
若 $X$ 为 $T_1$ 且 $x\in X$，则对每个 $y\ne x$，有含 $y$ 不含 $x$ 的开集 $U_y$；于是 $\{x\}^c=\bigcup_{y\ne x}U_y$ 开，故 $\{x\}$ 闭。反之，若 $\{x\}$ 闭，则 $\{x\}^c$ 是包含每个 $y\ne x$ 而不含 $x$ 的开集。
\end{proof}
% SUPPLEMENT: id=v14-ch04-t1-hierarchy-details; parent=ch04-t1-singletons; kind=proof+motivation+hypotheses; origin=independent-derivation
\begin{remark}
\textbf{动机与条件说明（重构新增）。}此判别把对每一对点选开集的要求，化成单点的闭性；不用可数性，所以把所有 $y\ne x$ 对应开集取并是合法且必要的一般步骤。它还表明 $T_1$ 空间中有限集闭，因为有限闭集并闭。$T_1$ 仍不保证 Hausdorff，余有限拓扑将给出反例。
\end{remark}
\begin{proof}
\textbf{分离层次的补写证明（重构新增）。}在正规空间中，给定闭集 $A$ 与 $x\notin A$，单点 $\{x\}$ 闭，故正规性可分离这两个闭集，从而正则。在正则空间中，对不同点 $x,y$，以闭集 $\{y\}$ 应用正则性，得互不相交的开邻域，故 Hausdorff。Hausdorff 空间中，用不交开邻域的其中一个区分有序对，得 $T_1$；$T_1$ 的开集又满足 $T_0$ 中的一个方向。因而有 $T_4\Rightarrow T_3\Rightarrow T_2\Rightarrow T_1\Rightarrow T_0$。
\end{proof}
% END SUPPLEMENT: id=v14-ch04-t1-hierarchy-details
% SOURCE p=132; printed=117; object=ch04-nonhausdorff-intro; region=middle
实际中出现的绝大多数拓扑空间（特别是本书中的空间）都是 Hausdorff 的，或经简单修改后成为 Hausdorff 空间。这里后一说法指的例子包括测度空间上的可积函数空间：把几乎处处相等的函数视为同一个元素后，$L^1$ 度量就使它成为 Hausdorff 空间。不过，两类通常不是 Hausdorff 的拓扑相当重要，值得特别提及：一类是等价类空间上的商拓扑，见\cref{section:ch04-continuous-maps}的\cref{exercise:ch04-quotient-topology,exercise:ch04-quotient-maps}；另一类是代数簇上的 Zariski 拓扑。这里不试图定义代数簇，而只描述向量空间上的 Zariski 拓扑。
% SOURCE p=132; printed=117; object=ch04-polynomial-maps; region=middle-lower
\begin{definition}[多项式映射]\label{definition:ch04-polynomial-map}
设 $k$ 为域，$k(X_1,\ldots,X_n)$ 为 $k$ 上 $n$ 个变量的多项式环。每个 $P\in k(X_1,\ldots,X_n)$ 通过用 $k$ 的元素代入形式未定元 $X_1,\ldots,X_n$，确定一个\textbf{多项式映射} $p:k^n\to k$。
\end{definition}
% SOURCE p=132; printed=117; object=ch04-polynomial-evaluation-injectivity; region=middle-lower
\begin{proposition}\label{proposition:ch04-polynomial-evaluation-injective}
对应 $P\mapsto p$ 是一一的，当且仅当 $k$ 无限。
\end{proposition}
% SUPPLEMENT: id=v14-ch04-polynomial-evaluation-details; parent=ch04-polynomial-evaluation-injectivity; kind=proof+motivation+hypotheses; origin=independent-derivation
\begin{remark}
\textbf{动机、符号与条件说明（重构新增）。}原文这里用圆括号 $k(X_1,\ldots,X_n)$ 表示多项式环；许多文献改用方括号表示该环，用圆括号表示有理函数域，阅读本段须依原文的明示约定。取 $n\ge1$；若没有变量，多项式只是常数，下面关于有限域的非单射结论不适用。形式多项式由系数确定，多项式映射由各点取值确定，二者未必相同；无限性正是使取值足够区分系数的条件。
\end{remark}
\begin{proof}
\textbf{补写证明（重构新增）。}一元非零多项式至多有其次数个根：若 $a$ 是根，用多项式除法写 $P=(X-a)Q$，在域中其余不同根也是 $Q$ 的根，再对次数归纳。若 $k$ 无限，一元非零多项式不能处处为零。对变量数归纳，将非零 $P$ 写为关于 $X_n$ 的多项式，其中某个系数 $C(X_1,\ldots,X_{n-1})$ 非零。由归纳假设，选前 $n-1$ 个坐标使 $C$ 的值非零；所得关于 $X_n$ 的一元多项式非零，再选末坐标使其值非零。因此非零多项式不能给出零映射，作用于两个多项式之差即得单射。

若 $k$ 有限，直接取非零形式多项式 $P(X_1,\ldots,X_n)=\prod_{a\in k}(X_1-a)$。它的最高次系数为 $1$，故不为零，却在每个 $k^n$ 点上都取值零。因此取值对应不单射。
\end{proof}
% END SUPPLEMENT: id=v14-ch04-polynomial-evaluation-details
% SOURCE p=132; printed=117; object=ch04-zariski-construction; region=lower
所有形如 $p^{-1}(\{0\})$ 的集合（$p$ 遍历全部多项式映射）组成的族对有限并封闭，因为
\begin{align*}
p^{-1}(\{0\})\cup q^{-1}(\{0\})=(pq)^{-1}(\{0\}),
\end{align*}
并且它包含 $k^n$（取 $p=0$）。因此由\cref{proposition:ch04-base-union-characterization,proposition:ch04-base-criterion}，所有形如 $\bigcap_{\alpha\in A}p_\alpha^{-1}(\{0\})$ 的集合（每个 $p_\alpha$ 为多项式映射），组成 $k^n$ 上某个拓扑的闭集族。
% SOURCE p=132; printed=117; object=ch04-zariski-definition; region=lower
\begin{definition}[Zariski 拓扑]\label{definition:ch04-zariski-topology}
上述以多项式映射族的共同零点集为闭集的拓扑，称为 $k^n$ 上的 \textbf{Zariski 拓扑}。
\end{definition}
% SOURCE p=132; printed=117; object=ch04-zariski-properties; region=lower
\begin{proposition}\label{proposition:ch04-zariski-properties}
Zariski 拓扑为 $T_1$。若 $k$ 有限，Zariski 拓扑为离散拓扑；若 $k$ 无限，Zariski 拓扑不是 Hausdorff 的，事实上，任意两个非空开集都有非空交。当 $n=1$ 时，Zariski 拓扑就是余有限拓扑。
\end{proposition}
% SOURCE p=132; printed=117; object=ch04-proof-zariski-properties; region=lower-bottom
\begin{proof}
由\cref{proposition:ch04-t1-singletons}，只需注意：若 $a=(a_1,\ldots,a_n)\in k^n$，则 $\{a\}=\bigcap_{j=1}^np_j^{-1}(\{0\})$，其中 $p_j(X_1,\ldots,X_n)=X_j-a_j$。无限域时非空开集相交的断言，只是多项式环 $k(X_1,\ldots,X_n)$ 为整环这一事实的另一种说法，即非零多项式 $P,Q$ 的积 $PQ$ 非零。
\end{proof}
% SUPPLEMENT: id=v14-ch04-zariski-details; parent=ch04-zariski-properties; kind=proof+motivation+hypotheses; origin=independent-derivation
\begin{proof}
\textbf{构造与省略结论的补写证明（重构新增）。}记 $Z(p)=p^{-1}(\{0\})$，$D(p)=k^n\setminus Z(p)$。域没有零因子，故 $p(x)q(x)=0$ 当且仅当至少一个因子为零，给出 $Z(pq)=Z(p)\cup Z(q)$，等价于 $D(pq)=D(p)\cap D(q)$。并且 $D(1)=k^n$，$D(0)=\varnothing$，所以 $\{D(p)\}$ 满足基判别。其所有并之补正是任意共同零点集，验证所宣称的闭集构造。

单点闭性由原证明已得。若 $k$ 有限，则 $k^n$ 有限，每个子集是有限个闭单点之并，故每个子集闭且开，拓扑离散。若 $k$ 无限，任意非空开集各含非空基成员 $D(p),D(q)$。对应形式多项式 $P,Q$ 非零，其积亦非零：可用多项式最高项的乘积系数非零，逐变量归纳验证整环性。由上一命题，$PQ$ 在某点取非零值，因此 $D(pq)\ne\varnothing$。所以原两个开集相交。$k^n$ 至少有两点，故这排除了 Hausdorff 性。

当 $n=1$，任意真闭集是某族零点集之交；若其中有非零形式多项式，交集包含于该多项式的有限根集；若全为零多项式，则交为全空间。另一方面任意有限集 $\{a_1,\ldots,a_r\}$ 是 $\prod_{j=1}^r(X-a_j)$ 的零点集，空集可用常数 $1$。所以闭集恰为有限集及全空间，正是余有限拓扑；在有限域上，这也等于离散拓扑。
\end{proof}
\begin{remark}
\textbf{动机与逐条件说明（重构新增）。}Zariski 拓扑把满足多项式方程作为“闭”的基本几何约束；方程组允许任意多条，故闭集取任意交。域的无零因子性质保证方程乘积描述两个零点集的并；若允许一般有零因子的环，这一步会失败，例如模 $6$ 的常数 $2,3$ 都无零点而乘积恒零。有限域与无限域的差别不是装饰性条件：前者的有限底空间与 $T_1$ 联合强制离散，后者由非零多项式不能处处为零而得到相反的非 Hausdorff 性。$n\ge1$ 同样用于保证无限域时空间不止一点。
\end{remark}
% END SUPPLEMENT: id=v14-ch04-zariski-details
% SOURCE p=132; printed=117; object=ch04-axioms-exercises-transition; region=bottom
习题中还有其他说明分离公理与可数性公理的例子。
% SOURCE p=132; printed=117; object=ch04-topology-exercises-title; region=bottom
\Needspace{14\baselineskip}
\subsection*{习题}
% SOURCE p=132; printed=117; object=ch04-ex01-t0-not-t1; number=1; region=bottom
\begin{exercise}\label{exercise:ch04-t0-not-t1}
若 $\operatorname{card}(X)\ge2$，则 $X$ 上存在一个 $T_0$ 而非 $T_1$ 的拓扑。
\end{exercise}
% SUPPLEMENT: id=v14-ch04-ex01-solution; parent=ch04-ex01-t0-not-t1; kind=solution; origin=independent-derivation
\begin{solution}
\textbf{补写解答（重构新增）。}固定 $p\in X$，令 $\mathcal T=\mathscr P(X\setminus\{p\})\cup\{X\}$。任意并若有成员为 $X$ 则为 $X$，否则不含 $p$；有限交若有成员不含 $p$ 则也不含 $p$，否则为 $X$。两端集合也在其中，故它是拓扑。若 $x\ne y$ 且 $y\ne p$，开集 $\{y\}$ 含 $y$ 不含 $x$；若 $y=p$，则 $x\ne p$，用开集 $\{x\}$ 从另一个方向区分。因此满足 $T_0$。取 $q\ne p$，唯一包含 $p$ 的开集为 $X$，它也包含 $q$，所以不能满足 $T_1$。至少两点用于选出 $q$；单点空间中唯一拓扑满足 $T_1$。
\end{solution}
% END SUPPLEMENT: id=v14-ch04-ex01-solution
% SOURCE p=133; printed=118; object=ch04-ex02-cofinite-countability; number=2; region=top
\begin{exercise}\label{exercise:ch04-cofinite-countability}
若 $X$ 为无限集合，则 $X$ 上的余有限拓扑为 $T_1$ 而非 $T_2$，并且它第一可数当且仅当 $X$ 可数。
\end{exercise}
% SUPPLEMENT: id=v14-ch04-ex02-solution; parent=ch04-ex02-cofinite-countability; kind=solution; origin=independent-derivation
\begin{solution}
\textbf{补写解答（重构新增）。}每个单点闭，故由\cref{proposition:ch04-t1-singletons}满足 $T_1$。任意两个非空开集的交之补有限，而 $X$ 无限，故其交非空；于是不同点不可能有不交开邻域，不满足 $T_2$。

若 $X$ 可数，对固定 $x$ 枚举 $X\setminus\{x\}=\{y_1,y_2,\ldots\}$，令 $U_n=X\setminus\{y_1,\ldots,y_n\}$。它们是 $x$ 的开邻域；任意开邻域 $V\ni x$ 的有限补集被某个有限初段包含，故某个 $U_n\subseteq V$，得到可数邻域基。

反之，若某点 $x$ 有可数邻域基 $\{U_n\}$，则 $C=\bigcup_nU_n^c$ 可数。对每个 $y\ne x$，开邻域 $X\setminus\{y\}$ 必包含某个 $U_n$，从而 $y\in U_n^c\subseteq C$。所以 $X\subseteq C\cup\{x\}$ 可数。事实上，只要一个点第一可数就已足够推出结论。
\end{solution}
% END SUPPLEMENT: id=v14-ch04-ex02-solution
% SOURCE p=133; printed=118; object=ch04-ex03-metric-normal; number=3; region=upper
\begin{exercise}\label{exercise:ch04-metric-normal}
每个度量空间都正规。（若 $A,B$ 为度量空间 $(X,\rho)$ 中的闭集，考虑满足 $\rho(x,A)<\rho(x,B)$ 或 $\rho(x,A)>\rho(x,B)$ 的点集。）
\end{exercise}
% SUPPLEMENT: id=v14-ch04-ex03-solution; parent=ch04-ex03-metric-normal; kind=solution; origin=independent-derivation
\begin{solution}
\textbf{补写解答（重构新增）。}度量空间单点闭：若 $y\ne x$，球 $B(y,\rho(x,y)/2)$ 不含 $x$，故单点之补开，从而空间为 $T_1$。现取不交闭集 $A,B$。若某个为空，分别用 $\varnothing,X$ 即可分离，所以设二者非空。

对非空集合 $A$，三角不等式取下确界给出
\begin{align*}
|\rho(x,A)-\rho(y,A)|\le\rho(x,y),
\end{align*}
故距离函数连续。闭性还给出 $\rho(x,A)=0$ 当且仅当 $x\in A$：若 $x\notin A$，闭集之补含一个以 $x$ 为中心的正半径球，故距离有正下界。因此
\begin{align*}
U&=\{x:\rho(x,A)<\rho(x,B)\}, &
V&=\{x:\rho(x,A)>\rho(x,B)\}
\end{align*}
是互不相交的开集，且 $A\subseteq U$、$B\subseteq V$。这验证正规性。分别使用每一点到另一闭集的正距离即可，不需要两个闭集之间存在统一的正距离。
\end{solution}
% END SUPPLEMENT: id=v14-ch04-ex03-solution
% SOURCE p=133; printed=118; object=ch04-ex04-rational-refinement; number=4; region=upper
\begin{exercise}\label{exercise:ch04-rational-refinement}
令 $X=\mathbb R$，并令 $\mathcal T$ 为所有形如 $U\cup(V\cap\mathbb Q)$ 的实数子集组成的族，其中 $U,V$ 是通常意义下的开集。则 $\mathcal T$ 是 Hausdorff 但不正则的拓扑。（结合\cref{exercise:ch04-metric-normal}，这表明比一个正规拓扑更强的拓扑，未必正规，甚至未必正则。）
\end{exercise}
% SUPPLEMENT: id=v14-ch04-ex04-solution; parent=ch04-ex04-rational-refinement; kind=solution; origin=independent-derivation
\begin{solution}
\textbf{补写解答（重构新增）。}把所有通常开集及所有形如 $V\cap\mathbb Q$ 的集合放在一起作为基候选。该族包含 $\mathbb R$，两个成员的交仍属该族，故满足基判别。其任意并可把通常开集部分合为 $U$、有理部分合为 $V\cap\mathbb Q$，所得拓扑恰为 $\mathcal T$。由于它包含通常拓扑，不同实数原来的不交开邻域仍开，所以它为 Hausdorff。

集合 $A=\mathbb R\setminus\mathbb Q$ 在新拓扑中闭，因为 $\mathbb Q$ 开。固定 $x\in\mathbb Q$。若 $W$ 是 $x$ 的新开邻域，写 $W=U\cup(V\cap\mathbb Q)$；无论 $x$ 落在 $U$ 还是 $V\cap\mathbb Q$，都能选通常开区间 $I\ni x$ 使 $I\cap\mathbb Q\subseteq W$。再设 $Z$ 是包含 $A$ 的新开集，写 $Z=U'\cup(V'\cap\mathbb Q)$。选无理数 $y\in I$；它必须在 $U'$ 中，故 $I\cap U'$ 是非空通常开集，包含某个有理数 $q$。于是 $q\in W\cap Z$。这说明 $x$ 与闭集 $A$ 无法用不交开集分离，故不正则。
\end{solution}
% END SUPPLEMENT: id=v14-ch04-ex04-solution
% SOURCE p=133; printed=118; object=ch04-ex05-separable-metric; number=5; region=upper-middle
\begin{exercise}\label{exercise:ch04-separable-metric-second-countable}
每个可分度量空间都第二可数。
\end{exercise}
% SUPPLEMENT: id=v14-ch04-ex05-solution; parent=ch04-ex05-separable-metric; kind=solution; origin=independent-derivation
\begin{solution}
\textbf{补写解答（重构新增）。}设 $D$ 为可数稠密集，令 $\mathcal B=\{B(d,q):d\in D,\ q\in\mathbb Q,\ q>0\}$，此族可数。给定开集 $U\ni x$，选 $r>0$ 使 $B(x,r)\subseteq U$，由稠密性选 $d\in D\cap B(x,r/3)$。记 $s=\rho(x,d)<r/3$，选有理数 $q$ 满足 $s<q<r-s$。于是 $x\in B(d,q)$；若 $y\in B(d,q)$，则 $\rho(x,y)\le s+\rho(d,y)<s+q<r$，故 $B(d,q)\subseteq U$。因此 $\mathcal B$ 为可数基。
\end{solution}
% END SUPPLEMENT: id=v14-ch04-ex05-solution
% SOURCE p=133; printed=118; object=ch04-ex06-upper-limit-topology; number=6; region=middle
\begin{exercise}\label{exercise:ch04-upper-limit-topology}
令 $\mathcal E=\{(a,b]:-\infty<a<b<\infty\}$。
\begin{enumerate}
\item\label{item:ch04-upper-limit-base} $\mathcal E$ 是 $\mathbb R$ 上某个拓扑 $\mathcal T$ 的基，且 $\mathcal E$ 的成员同时开闭。
\item\label{item:ch04-upper-limit-countability} $\mathcal T$ 第一可数但非第二可数。（若 $x\in\mathbb R$，则 $x$ 处每个邻域基都包含一个上确界为 $x$ 的集合。）
\item\label{item:ch04-upper-limit-density} $\mathbb Q$ 关于 $\mathcal T$ 在 $\mathbb R$ 中稠密。（所以\cref{proposition:ch04-second-countable-separable}的逆命题是错的。）
\end{enumerate}
\end{exercise}
% SUPPLEMENT: id=v14-ch04-ex06-solution; parent=ch04-ex06-upper-limit-topology; kind=solution; origin=independent-derivation
\begin{solution}
\textbf{补写解答（重构新增）。}
\begin{enumerate}
\item 每个 $x$ 属于 $(x-1,x]$，故该族覆盖。两个基候选的交若非空，等于 $(\max\{a,c\},\min\{b,d\}]$，仍在该族，故基判别成立。固定 $(a,b]$，对其补中的 $y$，若 $y\le a$，取任意 $(y-\varepsilon,y]$；若 $y>b$，取 $0<\varepsilon<y-b$。所得基邻域均包含于补集，故补集开，原集合闭。
\item $\{(x-1/n,x]:n\ge1\}$ 是 $x$ 处的邻域基：若 $x\in(a,b]$，选 $1/n<x-a$，则 $(x-1/n,x]\subseteq(a,b]$。若存在可数拓扑基 $\mathcal B$，对每个 $x$ 选 $B_x\in\mathcal B$ 满足 $x\in B_x\subseteq(x-1,x]$。于是 $\sup B_x=x$，不同 $x$ 对应不同 $B_x$，给出不可数集合 $\mathbb R$ 到可数族 $\mathcal B$ 的单射，矛盾。对任意邻域基在开邻域 $(x-1,x]$ 中细化，也证明了题目提示。
\Needspace{5\baselineskip}
\item 每个非空开集包含一个非退化 $(a,b]$，其内部区间 $(a,b)$ 含有理数，因此每个非空开集都遇到 $\mathbb Q$，即 $\mathbb Q$ 稠密。结合上一项，得到可分但非第二可数的空间。
\end{enumerate}
\end{solution}
% END SUPPLEMENT: id=v14-ch04-ex06-solution
% SOURCE p=133; printed=118; object=ch04-ex07-sequence-cluster; number=7; region=middle
\begin{exercise}\label{exercise:ch04-sequence-cluster-subsequence}
\begin{definition*}[序列的簇点]\phantomsection\label{definition:ch04-sequence-cluster-point}
设 $X$ 为拓扑空间。若对 $x$ 的每个邻域 $U$，都有无穷多个指标 $j$ 满足 $x_j\in U$，则称 $x\in X$ 为序列 $\{x_j\}$ 的一个\textbf{簇点}。
\end{definition*}
若 $X$ 第一可数，则 $x$ 为 $\{x_j\}$ 的簇点，当且仅当 $\{x_j\}$ 的某个子序列收敛于 $x$。
\end{exercise}
% SUPPLEMENT: id=v14-ch04-ex07-solution; parent=ch04-ex07-sequence-cluster; kind=solution+motivation+hypotheses; origin=independent-derivation
\begin{remark}
\textbf{定义的动机与条件说明（重构新增）。}簇点要求序列无限次回到每个邻域，而不要求最终留在其中；这里数的是指标，不是不同的取值。例如常值序列的值是簇点，但单点值域在 $T_1$ 空间中没有集合意义的聚点。第一可数性用于把无限次回访整理成一个真正收敛的子序列；逆方向不需此条件，正方向只需 $x$ 处有可数邻域基。
\end{remark}
\begin{solution}
\textbf{补写解答（重构新增）。}若 $x$ 为簇点，取递减邻域基 $\{U_n\}$。先选 $j_1$ 使 $x_{j_1}\in U_1$；已选 $j_n$ 时，由无穷多个指标的条件，选 $j_{n+1}>j_n$ 使 $x_{j_{n+1}}\in U_{n+1}$。对任意邻域 $V$，选 $U_N\subseteq V$，则所有 $n\ge N$ 都有 $x_{j_n}\in U_n\subseteq U_N\subseteq V$，所以子序列收敛于 $x$。反之，若某子序列收敛于 $x$，则每个邻域包含该子序列的所有充分后项；这些项有无穷多个不同指标，所以 $x$ 是原序列的簇点。
\end{solution}
% END SUPPLEMENT: id=v14-ch04-ex07-solution
% SOURCE p=133; printed=118; object=ch04-ex08-cofinite-sequence; number=8; region=middle-lower
\begin{exercise}\label{exercise:ch04-cofinite-sequence}
若 $X$ 是带余有限拓扑的无限集合，且 $\{x_j\}$ 是 $X$ 中两两不同的点组成的序列，则对每个 $x\in X$，都有 $x_j\to x$。
\end{exercise}
% SUPPLEMENT: id=v14-ch04-ex08-solution; parent=ch04-ex08-cofinite-sequence; kind=solution; origin=independent-derivation
\begin{solution}
\textbf{补写解答（重构新增）。}固定任意 $x\in X$ 及其邻域 $U$，选含 $x$ 的开集 $V\subseteq U$。$V$ 非空，所以 $X\setminus V$ 有限。由于序列的值两两不同，它落在 $X\setminus V$ 中的指标只有有限多个；取大于这些指标的 $J$（没有时任取 $J$），则所有 $j>J$ 都有 $x_j\in V\subseteq U$。故序列收敛于任意指定的 $x$。互异条件用来把有限个例外取值变成有限个例外指标；若常值为 $y\ne x$，开邻域 $X\setminus\{y\}$ 就阻止收敛于 $x$。
\end{solution}
% END SUPPLEMENT: id=v14-ch04-ex08-solution
% SOURCE p=133; printed=118; object=ch04-ex09-order-topology; number=9; region=lower
\begin{exercise}\label{exercise:ch04-order-topology}
\begin{definition*}[序拓扑]\phantomsection\label{definition:ch04-order-topology}
若 $X$ 为线性序集，由集合 $\{x:x<a\}$ 及 $\{x:x>a\}$（$a\in X$）生成的拓扑 $\mathcal T$ 称为\textbf{序拓扑}。
\end{definition*}
\begin{enumerate}
\item\label{item:ch04-order-separation} 若 $a,b\in X$ 且 $a<b$，则存在 $U,V\in\mathcal T$，使 $a\in U$、$b\in V$，且所有 $x\in U,y\in V$ 都满足 $x<y$。序拓扑是具有这一性质的最弱拓扑。
\item\label{item:ch04-order-relative} 若 $Y\subseteq X$，则 $Y$ 上的序拓扑不会强于由 $X$ 的序拓扑诱导的相对拓扑，但可以严格弱于后者。
\item\label{item:ch04-order-real} $\mathbb R$ 上的序拓扑是通常拓扑。
\end{enumerate}
\end{exercise}
% SUPPLEMENT: id=v14-ch04-ex09-solution; parent=ch04-ex09-order-topology; kind=solution+motivation+hypotheses; origin=independent-derivation
\begin{remark}
\textbf{定义的动机与条件说明（重构新增）。}序拓扑把“小于某个界”与“大于某个界”作为基本开条件；取生成拓扑自动补入有限交与任意并，也自动处理最大元、最小元及相邻点。线性序保证任意两点可比较，证明中的左右分离依赖这一点。子集的序仅记录剩余点之间的比较，而相对拓扑还保留原空间中被删去点形成的间隙，因此二者可能不同。
\end{remark}
\begin{solution}
\textbf{补写解答（重构新增）。}
\begin{enumerate}
\item 若存在 $c$ 满足 $a<c<b$，取 $U=\{x:x<c\}$、$V=\{y:y>c\}$。若不存在这样的 $c$，取 $U=\{x:x<b\}$、$V=\{y:y>a\}$；线性序及无中间点保证每个 $x\in U$ 都满足 $x\le a$，每个 $y\in V$ 都满足 $y\ge b$，故仍有 $x<y$。

若另一拓扑 $\mathcal S$ 有题述性质，对固定 $a$ 及每个 $x<a$，用该性质选 $\mathcal S$-开集 $U_x\ni x$ 和 $V_x\ni a$，使所有 $u\in U_x,v\in V_x$ 满足 $u<v$。代入 $v=a$ 得 $U_x\subseteq\{u:u<a\}$，故此下射线是这些 $U_x$ 的并，因而 $\mathcal S$-开。上射线同理也开。所以 $\mathcal S$ 包含全部生成子基，必须包含序拓扑，证明最弱性。
\item 对 $a\in Y$，$Y$ 中的下射线等于 $Y\cap\{x\in X:x<a\}$，上射线同理，故 $Y$ 的每个序子基成员在相对拓扑中开，由最小性得到包含关系。

严格包含的例子是通常有序的 $X=\mathbb R$ 和 $Y=\{0\}\cup(1,2)$。集合 $\{0\}=Y\cap(-1,1)$ 相对开。在 $Y$ 的序拓扑中，包含 $0$ 的子基射线只能是下射线，其端点 $y$ 位于 $(1,2)$；任何包含 $0$ 的有限子基交仍包含某个 $(1,y_0)$，因此不是单点。由子基构造定理，每个 $0$ 的序开邻域都含这样一个有限交，故 $\{0\}$ 不序开。
\item 实直线的上下射线通常开，所以序拓扑包含于通常拓扑。反之，通常拓扑的每个开区间 $(a,b)=\{x:x>a\}\cap\{x:x<b\}$ 在序拓扑中开；通常开集是这类开区间的并，故也序开。于是两个拓扑相等。
\end{enumerate}
\end{solution}
% END SUPPLEMENT: id=v14-ch04-ex09-solution
% SOURCE p=133,134; printed=118,119; object=ch04-ex10-connectedness; number=10; region=bottom-top
\begin{exercise}\label{exercise:ch04-connectedness}
\begin{definition*}[连通与不连通]\phantomsection\label{definition:ch04-connectedness}
若拓扑空间 $X$ 存在非空开集 $U,V$，使 $U\cap V=\varnothing$ 且 $U\cup V=X$，则称 $X$ \textbf{不连通}；否则称 $X$ \textbf{连通}。讨论 $X$ 的子集是否连通时，使用该子集的相对拓扑。
\end{definition*}
\begin{enumerate}
\item\label{item:ch04-connected-clopen} $X$ 连通当且仅当 $\varnothing,X$ 是 $X$ 中仅有的同时开闭的子集。
\item\label{item:ch04-connected-union} 若 $\{E_\alpha\}_{\alpha\in A}$ 是 $X$ 的连通子集族，且 $\bigcap_{\alpha\in A}E_\alpha\ne\varnothing$，则 $\bigcup_{\alpha\in A}E_\alpha$ 连通。
\item\label{item:ch04-connected-closure} 若 $A\subseteq X$ 连通，则 $\overline A$ 连通。
\item\label{item:ch04-connected-component} 每个点 $x\in X$ 都属于 $X$ 的唯一一个极大连通子集，且这个子集闭。
\end{enumerate}
\begin{definition*}[连通分支]\phantomsection\label{definition:ch04-connected-component}
上述包含 $x$ 的唯一极大连通子集称为 $x$ 的\textbf{连通分支}。
\end{definition*}
\end{exercise}
% SUPPLEMENT: id=v14-ch04-ex10-solution; parent=ch04-ex10-connectedness; kind=solution+motivation+hypotheses; origin=independent-derivation
\begin{remark}
\textbf{定义的动机与逐条件说明（重构新增）。}连通性禁止把空间分成两块彼此隔开的开放部分。“非空”排除 $\varnothing$ 与全空间这类平凡分割；“不交”保证确实分成两块；“并为全空间”要求分割覆盖全部空间，而不只描述两个互不相交的小区域。对子集必须用相对拓扑，否则在环境中不开的子集会被误判。连通分支用包含关系的极大性，把不能再扩大的连通部分作为空间的组成单位；“极大”不表示某种数值大小最大。
\end{remark}
\begin{solution}
\textbf{补写解答（重构新增）。}
\begin{enumerate}
\item 若 $X=U\cup V$ 是题述分割，则 $U$ 非空、非全空间，且其补 $V$ 开，故 $U$ 同时开闭。反之，若存在非空真子集 $U$ 同时开闭，则 $U$ 与 $X\setminus U$ 是这样的分割。
\item 若指标族为空，则并为空空间，它不存在分成两非空开集的分割，故结论成立。否则取公共点 $p\in\bigcap_\alpha E_\alpha$，记 $E=\bigcup_\alpha E_\alpha$。若 $E$ 有相对开分割 $G,H$，可设 $p\in G$。对每个 $\alpha$，$G\cap E_\alpha,H\cap E_\alpha$ 是 $E_\alpha$ 中相对开、不交且覆盖它的两集，前者非空；连通性迫使后者为空。故 $H$ 不含任何 $E_\alpha$ 的点，即 $H=\varnothing$，矛盾。
\item 若 $\overline A$ 有非空相对开分割 $G,H$，则二者都遇到 $A$。例如写 $G=O\cap\overline A$，其中 $O$ 在 $X$ 中开；取 $x\in G$，因 $x\in\overline A$，邻域 $O$ 遇到 $A$，故 $G\cap A\ne\varnothing$。同理 $H\cap A\ne\varnothing$。于是二者与 $A$ 的交给出 $A$ 的相对开分割，矛盾。$A=\varnothing$ 时结论直接成立。
\item 令 $C_x$ 为全部包含 $x$ 的连通子集之并。候选族非空，因为 $\{x\}$ 连通；这些集合共有点 $x$，由前述并集结论，$C_x$ 连通。任何包含 $x$ 的连通子集都已包含在 $C_x$ 中，所以 $C_x$ 为唯一极大者。又由闭包结论，$\overline{C_x}$ 连通且含 $x$，故 $\overline{C_x}\subseteq C_x$，从而 $C_x$ 闭。
\end{enumerate}
\end{solution}
\begin{remark}
\textbf{结论的条件说明（重构新增）。}并集结论中的共同点防止各连通部分彼此分开；删掉它，例如实直线的两个不同单点之并不连通。共同交非空是充分条件，未宣称必要。闭包结论利用 $A$ 在其闭包中稠密，不需要 $T_1$ 或 Hausdorff。各连通分支还两两相同或不交：若两分支相交，其并由并集结论连通，极大性使两者相等。
\end{remark}
% END SUPPLEMENT: id=v14-ch04-ex10-solution
% SOURCE p=134; printed=119; object=ch04-ex11-finite-union-closure; number=11; region=upper
\begin{exercise}\label{exercise:ch04-finite-union-closure}
若 $E_1,\ldots,E_n$ 为拓扑空间的子集，则 $\bigcup_{j=1}^nE_j$ 的闭包为 $\bigcup_{j=1}^n\overline{E_j}$。
\end{exercise}
% SUPPLEMENT: id=v14-ch04-ex11-solution; parent=ch04-ex11-finite-union-closure; kind=solution; origin=independent-derivation
\begin{solution}
\textbf{补写解答（重构新增）。}每个 $E_j$ 包含于 $E=\bigcup_iE_i$，所以 $\overline{E_j}\subseteq\overline E$；取并得到一向包含。反之，$F=\bigcup_j\overline{E_j}$ 是有限个闭集之并，故闭，且含 $E$，由闭包的最小性，$\overline E\subseteq F$。有限性用于保证该并仍闭，不能任意推广到可数并：在 $\mathbb R$ 中取 $E_n=\{1/n\}$，并集的闭包含 $0$，各闭包的并却不含 $0$。
\end{solution}
% END SUPPLEMENT: id=v14-ch04-ex11-solution
% SOURCE p=134; printed=119; object=ch04-ex12-kuratowski; number=12; region=upper-middle
\begin{exercise}\label{exercise:ch04-kuratowski-closure}
\begin{definition*}[Kuratowski 闭包算子]\phantomsection\label{definition:ch04-kuratowski-closure}
设 $X$ 为集合。$X$ 上的一个 \textbf{Kuratowski 闭包算子}是映射 $A\mapsto A^*$，从 $\mathscr P(X)$ 到自身，满足：
\begin{enumerate}[label=(\roman*)]
\item $\varnothing^*=\varnothing$；
\item 对所有 $A$，有 $A\subseteq A^*$；
\item 对所有 $A$，有 $(A^*)^*=A^*$；
\item 对所有 $A,B$，有 $(A\cup B)^*=A^*\cup B^*$。
\end{enumerate}
\end{definition*}
\begin{enumerate}
\item\label{item:ch04-kuratowski-from-topology} 若 $X$ 为拓扑空间，则映射 $A\mapsto\overline A$ 是 Kuratowski 闭包算子。（使用\cref{exercise:ch04-finite-union-closure}。）
\item\label{item:ch04-kuratowski-to-topology} 反之，给定 Kuratowski 闭包算子，令 $\mathcal F=\{A\subseteq X:A=A^*\}$ 及 $\mathcal T=\{U\subseteq X:U^c\in\mathcal F\}$。则 $\mathcal T$ 是一个拓扑，并且对任意 $A\subseteq X$，$A^*$ 正是它关于 $\mathcal T$ 的闭包。
\end{enumerate}
\end{exercise}
% SUPPLEMENT: id=v14-ch04-ex12-solution; parent=ch04-ex12-kuratowski; kind=solution+motivation+hypotheses; origin=independent-derivation
\begin{remark}
\textbf{定义的动机与逐条件说明（重构新增）。}四条公理允许把“取闭包”作为原始操作来重建拓扑。保持空集排除无缘无故产生必在的点；扩张性要求保留原集合；幂等性表示补成闭集后再作一次不会增加点；有限并保持性对应闭集对有限并封闭。映射值仍在 $\mathscr P(X)$ 中固定了底空间，也与扩张性共同给出 $X^*=X$。这里不预设单调性，它将由有限并保持性推出，因而无需多加一条公理。
\end{remark}
\begin{solution}
\textbf{补写解答（重构新增）。}
\begin{enumerate}
\item 空集本身闭，故其闭包为空。定义直接给出 $A\subseteq\overline A$；$\overline A$ 已经闭，所以再取闭包不变。最后，\cref{exercise:ch04-finite-union-closure}在 $n=2$ 时给出有限并保持性。
\item 先证单调性。若 $A\subseteq B$，则 $B=A\cup B$，因而
\begin{align*}
B^*=(A\cup B)^*=A^*\cup B^*,
\end{align*}
所以 $A^*\subseteq B^*$。由空集公理及扩张性有 $\varnothing,X\in\mathcal F$。若 $F,G\in\mathcal F$，则 $(F\cup G)^*=F^*\cup G^*=F\cup G$，故有限并保持在 $\mathcal F$ 中。对非空指标族 $\{F_\alpha\}\subseteq\mathcal F$，记 $H=\bigcap_\alpha F_\alpha$；单调性给出 $H^*\subseteq F_\alpha^*=F_\alpha$ 对每个 $\alpha$ 成立，故 $H^*\subseteq H$，再用扩张性得等号。空交是已属于 $\mathcal F$ 的 $X$。因此 $\mathcal F$ 满足闭集公理，由 de Morgan 律，其补集族 $\mathcal T$ 满足拓扑公理。

幂等性说明 $A^*\in\mathcal F$，扩张性说明它含 $A$。若 $F\in\mathcal F$ 且 $A\subseteq F$，单调性又给出 $A^*\subseteq F^*=F$。所以 $A^*$ 是包含 $A$ 的最小闭集，恰为新拓扑中的闭包。
\end{enumerate}
\end{solution}
% END SUPPLEMENT: id=v14-ch04-ex12-solution
% SOURCE p=134; printed=119; object=ch04-ex13-dense-open-closure; number=13; region=middle
\begin{exercise}\label{exercise:ch04-dense-open-closure}
若 $X$ 为拓扑空间，$U$ 在 $X$ 中开，且 $A$ 在 $X$ 中稠密，则 $\overline U=\overline{U\cap A}$。
\end{exercise}
% SUPPLEMENT: id=v14-ch04-ex13-solution; parent=ch04-ex13-dense-open-closure; kind=solution; origin=independent-derivation
\begin{solution}
\textbf{补写解答（重构新增）。}$U\cap A\subseteq U$ 立即给出 $\overline{U\cap A}\subseteq\overline U$。对 $x\in U$，任意开邻域 $V\ni x$ 与 $U$ 的交非空且开，由 $A$ 稠密，有 $V\cap U\cap A\ne\varnothing$。故 $x\in\overline{U\cap A}$，即 $U\subseteq\overline{U\cap A}$。右端闭，取闭包便得反向包含。

开性用于保证 $V\cap U$ 可作为稠密性的测试集；若删去，在 $\mathbb R$ 中取 $A=\mathbb Q$、$U=\{\sqrt2\}$，则右端为空而左端非空。稠密性也不可无条件删除，例如 $A=\varnothing$、$U=X\ne\varnothing$ 即失败。
\end{solution}
% END SUPPLEMENT: id=v14-ch04-ex13-solution
\end{document}
